\chapter{资源配置、优化与处方决策}
\label{ch:week09-optimization}

企业在选择行动时受到预算、设备、人力、库存和交期等条件约束。多个行动可能
竞争同一资源，因而需要比较整体配置方案。预测与因果分析提供需求、风险和效果
估计，优化模型根据目标函数与资源条件确定决策变量的取值。

线性规划对偶给出资源价格与最优性证明；整数约束处理离散选择和活动
先后关系。随机库存模型与尾部风险度量用于不确定环境下的方案比较，决策导向
学习则以行动损失评价估计参数。求解结果还需经过可行性检查、参数敏感性分析
与实施条件核验。

\section{管理资源的优化配置}

预测模型输出未来状态，例如某 SKU 下周需求为 30 件，或某客户流失概率为 0.8。因果推断评估行动效果，例如发放优惠券是否提高留存。优化根据预算、库存、产能、人力、服务水平和风险约束选择行动。

一个抽象的管理决策可以写为
\[
  a^*(x)
  =
  \argmin_{a\in\mathcal A(x)}
  \left\{\E[L(Y(a),a)\mid X=x]+\lambda C(a)\right\},
\]
其中 \(X=x\) 是当前信息，\(Y(a)\) 是采取行动 \(a\) 后的潜在结果，\(\mathcal A(x)\) 是可行行动集合。\(C(a)\) 只计入尚未包含在 \(L\) 中的成本，避免重复扣减。若需求等外生状态不受行动影响，可将随机结果写成该外生状态与行动共同决定的损失；若行动会改变需求，则需要第~\ref{ch:week08-causality}~章的干预分布，并根据该分布评价行动后的结果。

\begin{definition}[处方分析]
处方分析是在给定数据、模型、成本、收益和约束的条件下，给出可执行行动建议的分析过程。其输出包括行动选择、可行性依据、资源使用和相应风险。
\end{definition}

处方分析的输出可以是补货量、排班计划、产能分配、价格方案、客户触达清单、供应商选择或配送路线。其共同特征是：行动集合受到约束，目标函数能够明确或近似表达，且建议需要由组织落实。行动变量可以连续取值，并不要求候选方案只有有限个。

\section{可行域、目标与经济解释}
\subsection{优化模型的基本结构}
\subjectindex{kexingyu}{可行域}
\subjectindex{mubiaohanshu}{目标函数}

一个标准优化问题可以写作：
\[
  \max_{x} f(x)
\]
约束为
\[
  g_j(x)\le b_j,\qquad j=1,\ldots,m,
\]
\[
  x\in\mathcal X.
\]
这里 \(x\) 是决策变量，\(f(x)\) 是目标函数，约束 \(g_j(x)\le b_j\) 定义可行域，\(\mathcal X\) 给出变量类型，例如连续、整数或二元。

\begin{definition}[可行解与最优解]
满足所有约束的 \(x\) 称为可行解。若可行解 \(x^*\) 使目标函数在可行域中达到最大或最小，则称为最优解。
\end{definition}

线性规划要求目标函数和约束都是线性的：
\[
  \max_x c^\top x
  \quad
  \text{s.t.}
  \quad
  Ax\le b,\quad x\ge 0.
\]
若部分变量必须取整数或 0-1 值，则得到整数规划或 0-1 规划。企业决策中常见的排班、选址、补货批量、配送路径和客户触达名额往往具有整数约束。整数约束使问题更贴近业务，也通常使求解更困难。

\subsection{目标函数与成本核算}

目标函数表示组织希望改善的结果。短期收入、利润、服务水平和长期客户价值对应不同目标；选择目标后，还需明确成本归属与时间范围。其他必须达到的要求可写为约束。

库存补货的一个目标可以写为期望净价值：
\[
  \E[V(q,D)]
  =
  m\E[\min(I+q,D)]
  -
  h\E[(I+q-D)^+]
  -
  p\E[(D-I-q)^+]
  -
  cq.
\]
其中 \(q\) 是订货量，\(D\) 是随机需求，\(I\) 是当前库存，\(x^+=\max(x,0)\)。
\(m\) 表示扣除其他可变销售成本、尚未扣除本次采购成本的单位销售收入，
\(h\) 是单位持有成本，\(p\) 是单位缺货惩罚，\(c\) 是单位采购成本。
因此本次采购支出单独计为 \(cq\)，已有库存的购入支出视为沉没成本。
本例取期末残值为0。配套数据字段 \texttt{contribution\_margin} 按这里的 \(m\) 定义使用。

\begin{modernbox}[管理解释]
同一需求分布在不同成本条件下，对应的最优补货量可能不同。提高缺货惩罚会增加备货收益，提高持有成本会增加积压代价，库容则限定可行数量。目标函数和约束共同决定方案。
\end{modernbox}

目标函数还应区分收益和现金流。贡献毛利、采购现金支出、库存占用、缺货惩罚和客户体验损失可能发生在不同时间。若企业现金流紧张，预算约束可能比期望利润更重要；若客户服务水平是战略目标，缺货惩罚应反映长期客户价值，而不只是单次订单损失。

\subsection{约束与可执行性}

企业优化问题的约束描述资源限制、服务要求和业务规则。常见约束包括：
\begin{enumerate}
  \item 预算约束：采购、营销、人工或运输预算不能超限；
  \item 产能约束：机器、人力、仓库、客服或配送能力有限；
  \item 服务约束：缺货率、准时率、响应时间或最低覆盖率必须达标；
  \item 业务规则：最小起订量、整数包装、渠道优先级或合同承诺；
  \item 风险约束：单一供应商暴露、极端库存、客户公平性或合规要求。
\end{enumerate}

约束决定方案可行性。库存算例包括预算和库容限制；应用于具体采购问题时，可按合同和供给条件加入提前期、供应商产能、最小批量及产品替代关系。

\section{线性规划与资源的影子价格}
\subjectindex{yingzijiage}{影子价格}
\subjectindex{xianxingguihua}{线性规划}

\subsection{资源配置的原问题}
某工厂安排两类任务，数量为 \(x_1,x_2\ge0\)，收益分别为 3、2。两类资源的限制为 \(x_1+x_2\le4\) 和 \(2x_1+x_2\le5\)。教学模型是
\[
\max 3x_1+2x_2
\quad\text{s.t.}\quad
x_1+x_2\le4,\quad2x_1+x_2\le5,\quad x\ge0.
\]
两个约束交点为 \((1,3)\)，目标值为 9。该可行解的最优性还需由对偶上界验证。

\subsection{弱对偶与最优性证明}
\subjectindex{duiou}{对偶}
\subjectindex{zuiyouxingzhengming}{最优性证明}
对一般原问题 \(\max\{c^\top x:Ax\le b,x\ge0\}\)，选取非负资源价格 \(y\ge0\)，使 \(A^\top y\ge c\)。那么任一可行 \(x\) 都满足
\[
c^\top x\le y^\top Ax\le y^\top b.
\]
因此最紧的这类上界由对偶问题
\[
\min b^\top y\quad\text{s.t.}\quad A^\top y\ge c,\ y\ge0
\]
给出。本例取 \(y=(1,1)\)，两个产品的不等式都取等号，上界为 \(4+5=9\)，与已找到的可行收益一致，证明 \((1,3)\) 最优。在线性规划可行且最优值有限等标准条件下，强对偶保证最优原、对偶值相等\cite{bertsimas1997}。

\subsection{影子价格的有效区间}
将第一类资源上限改为 \(4+\Delta\)，两条约束交点为
\(x_1=1-\Delta,x_2=3+2\Delta\)。当 \(0\le\Delta\le1\) 时，该点可行，
收益为 \(9+\Delta\)，仍等于对偶价格 \(y=(1,1)\) 给出的上界。因此这一范围内
资源增量价值为每单位1。超过该范围后需重新求解。整数约束下的价值函数可以
呈阶梯状，其有限增量与连续模型的局部导数应分别计算。

\begin{figure}[htbp]
\centering
\begin{tikzpicture}
\begin{axis}[width=0.84\linewidth,height=7cm,xlabel={任务1数量 \(x_1\)},ylabel={任务2数量 \(x_2\)},
xmin=0,xmax=3,ymin=0,ymax=5,xtick={0,1,2,2.5,3},
legend style={at={(1.02,1)},anchor=north west}]
\addplot[fill=BDMTeal!10,draw=BDMTeal] coordinates {(0,0)(2.5,0)(1,3)(0,4)(0,0)};
\addlegendentry{可行域}
\addplot[BDMAmber,dashed,domain=0:3] {4.5-1.5*x};
\addlegendentry{目标值 9}
\addplot[only marks,mark=*,BDMAmber,mark size=2.7pt,forget plot] coordinates {(1,3)};
\node[font=\small,anchor=south west,fill=white,inner sep=3pt]
  at (axis cs:1.05,3.1) {最优点 \((1,3)\)};
\node[font=\footnotesize,text=BDMTeal,anchor=center]
  at (axis cs:.85,1.2) {收益 \(3x_1+2x_2\le9\)};
\end{axis}
\end{tikzpicture}
\caption{最优点与支持它的等目标线；资源价格提供独立的上界证明。}
\end{figure}

\subsection{边际价值与资源分配}

边际价值衡量增加一单位资源带来的目标函数改善，可用于比较不同用途的增量收益。

对 SKU \(k\)，设订货量为 \(q\)，期望净价值为 \(V_k(q)\)。第 \(q\) 个订货单位的边际价值为
\[
  \Delta V_k(q)=V_k(q)-V_k(q-1).
\]
在没有联合约束时，可以为每个 SKU 单独选择最大化 \(V_k(q)\) 的订货量；有预算和库容约束时，则需比较各 SKU 占用共同资源的收益。边际价值高但占用库容大的产品不一定总是优先，因为约束维度不止一个。

边际价值还可以帮助解释优化结果。若某个产品前几个单位的边际价值很高，而后续单位的边际价值迅速下降，继续增加订货量所带来的增量收益就会逐渐降低；若另一个产品边际价值较低但占用预算和库容很少，也可能成为约束紧张时的优先选择。

\subsection{离散案例中的约束价值}

对于本章的收益最大化问题，在目标和其他条件不变时，收紧资源约束不会提高最优值，
放松约束不会降低最优值；若被改变的约束并不限制当前方案，最优值也可能不变。
线性规划的对偶价格刻画局部资源价值，小规模离散案例则可通过有限增量比较。

\begin{definition}[资源增量价值]
若某约束右端从 \(b\) 放宽为 \(b+\Delta b\)，最优目标值从 \(z(b)\) 变为 \(z(b+\Delta b)\)，则在 \(\Delta b>0\) 下，单位资源增量对应的平均目标改善为
\[
  \frac{z(b+\Delta b)-z(b)}{\Delta b}.
\]
若价值函数在该点可微，令 \(\Delta b\to0\) 可得到局部边际价值。
离散问题的差商一般依赖增量大小，不能无条件等同于线性规划的对偶影子价格。
\end{definition}

影子价格反映约束资源在适用范围内的增量价值。若预算约束的影子价格很高，说明额外采购预算可能显著提高期望净价值；若库容约束的影子价格很高，说明扩容、临时仓或调整 SKU 结构可能值得讨论。反之，如果放松某约束几乎不改变最优值，该资源在所考察范围内的增量价值较低。

\section{随机需求、尾部风险与稳健性}
\subsection{确定性、随机性与情景}

需求、到货时间、客户响应率和产能可用性都具有不确定性。以平均值替代其分布可能低估高损失情景的影响。例如平均需求为 20 的产品，在缺货惩罚很高时，按均值订货可能产生较高缺货损失。

库存案例采用离散情景表示需求不确定性。设需求有 \(S\) 个情景，情景 \(s\) 的概率为 \(\pi_s\)，需求为 \(d_s\)，则订货量 \(q\) 的期望净价值为
\[
  \sum_{s=1}^{S}\pi_s V(q,d_s).
\]
其中 \(\pi_s\ge0\) 且 \(\sum_s\pi_s=1\)。情景可以来自历史分布、预测模型或压力测试；构造情景时需覆盖对成本与服务水平有重要影响的需求和供给变化。

\subsection{报童模型的临界分位数}
\subjectindex{baotongmoxing}{报童模型}
先研究没有初始库存、一次订货 \(q\)、随机需求 \(D\) 的问题。设售价为 \(p_s\)，采购成本为 \(c\)，剩余残值为 \(v\)，且 \(v<c<p_s\)。期望利润为
\[
\Pi(q)=p_s\E[\min(q,D)]+v\E[(q-D)^+]-cq.
\]
多订一单位，在需求超过库存时赚 \(p_s-c\)，在需求不足时损失 \(c-v\)。若需求分布连续且可微分处理合法，
\[
\Pi'(q)=(p_s-c)(1-F_D(q))-(c-v)F_D(q).
\]
令导数为零得到
\[
F_D(q^*)=\frac{p_s-c}{p_s-v}.
\]
售价 10、采购成本 6、残值 2 时，临界分位数为 \(0.5\)；若另有未满足需求的单位惩罚 4，缺货损失增加到 8、过量损失仍为 4，临界分位数为 \(2/3\)。最优订货由分布分位数与经济成本共同决定，通常不等于平均需求。离散需求使用满足累计概率条件的分位数，并检查整数边界。

\subsection{条件风险价值}
\subjectindex{CVaR}{CVaR（条件风险价值）}
两个方案可能有相同期望收益，却承担不同极端风险。对可积\footnote{随机变量可积是指其绝对值的期望有限，因而其期望损失具有有限、明确的数值。}损失随机变量 \(L\) 和置信水平 \(\alpha\in(0,1)\)，定义
\[
\operatorname{CVaR}_\alpha(L)
=\min_{\eta\in\R}
\left\{\eta+\frac1{1-\alpha}\E[(L-\eta)^+]\right\}.
\]
这一定义也适用于分位点具有正概率质量的离散分布。\footnote{风险价值（value at risk，VaR）是损失分布在指定置信水平下的分位数；条件风险价值（conditional value at risk，CVaR）综合最坏一部分概率质量的损失。离散分布在分位点可能集中一段概率质量，因此须按定义处理边界。}因此，CVaR 不能一概写成条件期望 \(\E[L\mid L\ge\operatorname{VaR}_\alpha]\)。对有限情景，引入 \(u_s\ge L_s-\eta,u_s\ge0\)，即可把超额项转为线性约束\cite{cvar09}。最小化期望损失加 \(\lambda\operatorname{CVaR}\) 表达一种风险偏好；\(\lambda\) 的选择应由管理者解释。

\begin{example}[分位点有概率质量时]
三个等概率情景的损失为 \(0,0,10\)，取 \(\alpha=0.5\)。\(\operatorname{VaR}_{0.5}=0\)，但 \(\E[L\mid L\ge0]=10/3\) 包含了全部概率质量。变分目标在 \(\eta=0\) 处取得最小值 \(0+2\times10/3=20/3\)。这一数值对应最坏的一半概率质量：损失 10 的全部 \(1/3\)，再加损失 0 中的 \(1/6\)。这正是离散分布下不能机械使用尾部条件均值的原因。
\end{example}

\subsection{分布稳健优化}
\subjectindex{fenbulubang}{分布鲁棒优化}
情景概率本身也可能估错。分布鲁棒优化把一个估计分布扩展成集合 \(\mathcal U\)，求
\[
\min_{a\in\mathcal A}\sup_{Q\in\mathcal U}\E_Q[L(a,\xi)].
\]
集合可以通过矩约束\footnote{矩约束限定均值、方差等分布特征；距离球则收集与参考分布的指定距离不超过给定半径的分布。二者都是表达分布不确定范围的方式。}、距离球或业务压力情景构造。集合越大通常越保守；过小可能漏掉风险，过大可能让任何行动都显得不可接受。有限样本性质取决于分布集合和半径的构造，应按所用方法的条件选择并报告这些参数。

\subsection{敏感性分析与鲁棒性}
\subjectindex{minganxing}{敏感性分析}

需求预测、成本参数、缺货惩罚、预算和库容都可能变化。敏感性分析考察关键参数变化时最优方案的稳定性。

敏感性分析可检查以下变化：
\begin{itemize}
  \item 预算上限降低时，哪些 SKU 被压缩；
  \item 库容上限降低时，是否转向体积更小或边际价值更高的 SKU；
  \item 缺货惩罚提高时，模型是否增加安全库存；
  \item 需求高情景概率提高时，补货量是否显著变化；
  \item 单个参数微小变化是否导致方案大幅跳变。
\end{itemize}

鲁棒性关注方案在不利情景下是否仍可接受。一个期望值最高的方案可能在高需求或供应延迟情景下损失很大。管理者可以采用保守情景、预算缓冲、服务水平约束或最大后悔值\footnote{某情景下的后悔值，是方案损失与事先知道该情景时可达的最小损失之差。最小最大后悔准则选择最坏情景后悔值最小的方案。}控制风险。鲁棒方案不一定有最高期望收益，但可能更适合执行风险较高的环境。

\section{制造排程与决策导向学习}

\subsection{将先后顺序写成整数约束}
\subjectindex{MILP}{MILP（混合整数线性规划）}
\subjectindex{paicheng}{排程}
混合整数线性规划（MILP）\footnote{MILP 是 mixed-integer linear programming 的缩写，指目标和约束为线性、部分变量取整数而其余变量连续的优化模型，常用于选择、排班和生产排程。参见 \url{https://developers.google.com/optimization/mip/mip_example}。}可用二元变量表达任务的先后关系。考虑两项任务 A、B 共用一台设备，处理时间 \(p_A=3,p_B=2\)，交期 \(d_A=3,d_B=4\)，延期权重 \(w_A=1,w_B=3\)。令开始时间为 \(S_i\)，完成时间 \(C_i=S_i+p_i\)，延期量 \(T_i\ge C_i-d_i,T_i\ge0\)。二元变量 \(z=1\) 表示 A 在 B 前，写成
\[
S_B\ge C_A-M(1-z),\qquad
S_A\ge C_B-Mz.
\]
另要求 \(S_i\ge0,C_i\le H,z\in\{0,1\}\)。在本例 \(H=5\) 且可行，取 \(M=H\) 足以使未启用的约束失效；一般问题需要根据有效时界推导 \(M\)，任意取巨大数字会损害数值稳定性。

目标为 \(\min(T_A+3T_B)\)。A 先执行时完成于 3、5，目标为 3；B 先执行时完成于 2、5，目标为 2。因此最早交期优先在加权总延期目标下不保证最优。增加设备资格、维护窗口、工艺前序与换线成本后，模型会扩展为实际排程问题。该算例未计入这些附加约束。

\subsection{预测损失与决策损失}
\subjectindex{juecedaoxiangxuexi}{决策导向学习}
\subjectindex{juecehouhui}{决策后悔}
设成本向量为 \(c\)，可行决策为 \(x\in\mathcal X\)，根据预测 \(\hat c\) 选择 \(x^*(\hat c)\)。决策后悔值为
\[
\ell_{\rm dec}(\hat c,c)
=c^\top x^*(\hat c)-c^\top x^*(c)\ge0.
\]
若两个互斥行动真实成本为 \((1,1.1)\)，预测 \((1.06,1.04)\) 的误差很小，却选择第二项，后悔值为 0.1；预测 \((2,2.2)\) 误差更大，却保持正确排序、后悔为零。两种预测的误差与决策后悔值排序不同。

Smart Predict-then-Optimize 将优化结构引入学习损失\cite{spo09}；数据驱动处方分析研究如何从协变量学习决策\cite{bk2020}；2025 年 DF2 等工作继续探索决策导向学习的计算与分布问题\cite{df209}。这些方法的核心挑战是离散解映射的不连续、可微近似、训练成本和分布漂移。实证比较应在相同数据与决策任务下，评价联合学习方法，以及先训练预测模型、再调用求解器的方法。

\subsection{滚动优化与仿真}
\subjectindex{gundongyouhua}{滚动优化}
Oracle 与 Red Bull 的公开案例将仿真用于赛前与赛中策略响应\cite{oracle01}。制造排程可采用滚动时域：每次用最新订单和设备状态求未来 \(H\) 小时计划，只执行近期承诺的部分，下一时刻再更新。为避免频繁改计划扰乱现场，可加入计划变更成本或冻结窗口\footnote{冻结窗口是计划期内暂不允许常规调整的近期时段，用于保护已经投入准备或执行的任务；紧急变更应另设审批规则。}。

仿真器应保存随机种子和情景版本，并检查对故障、需求与换线的结构假设。仿真最优解依赖模拟模型的结构和参数，其实际效果需经现场验证。优化得出的收益也应与人工基准或现行规则比较，并报告相同情景下的成本或收益差额。

\section{延期风险与加急组合选择}
\label{sec:expedite-choice}
\subjectindex{zengliangshouyi}{增量收益}
订单加急组合选择同时需要风险预测、干预效果和行动约束。设工厂需要在四张订单中
选择至多两张加急。\(p_i\) 表示不加急时的延期概率，\(\delta_i\) 表示加急使延期
概率下降的绝对幅度，\(v_i\) 为一次延期的损失，\(c_i\) 为加急成本。
以下为给定风险、效果与成本参数的算例，金额单位为千元。各 \(\delta_i\) 满足
\(0\le\delta_i\le p_i\)，其估计依据是第8章的处理效应分析。

\begin{table}[htbp]
\centering
\caption{四张订单的风险、加急效果与增量净收益}
\label{tab:expedite-choice}
\begin{tabular}{lrrrrr}
\toprule
\rowcolor{BDMTableHead}
订单 & \(p_i\) & \(\delta_i\) & \(v_i\) & \(c_i\) & \(v_i\delta_i-c_i\)\\
\midrule
A & 0.60 & 0.10 & 10 & 0.40 & 0.60\\
B & 0.40 & 0.25 & 8 & 0.60 & 1.40\\
C & 0.30 & 0.20 & 5 & 0.30 & 0.70\\
D & 0.10 & 0.05 & 4 & 0.20 & 0.00\\
\bottomrule
\end{tabular}
\end{table}

不加急的期望损失是 \(v_ip_i\)，加急后是 \(v_i(p_i-\delta_i)+c_i\)。两者相减，
得到增量净收益 \(b_i=v_i\delta_i-c_i\)。风险 \(p_i\) 在这一简化的净收益差值中抵消，
因此，在本例固定损失与可加效果的假设下，增量收益取决于延期概率的可降低幅度。
令 \(z_i\in\{0,1\}\) 表示是否加急，问题成为
\begin{equation}
\max_z\ \sum_{i\in\{A,B,C,D\}}b_iz_i,
\qquad \sum_i z_i\le2.
\label{eq:expedite-knapsack}
\end{equation}
按风险最高选择 A、B，收益为2.00；比较净收益则选择 B、C，收益为2.10。
风险预测见第~\ref{ch:week07-prediction}~章，效果识别见第~\ref{ch:week08-causality}~章。
组合优化将这些估计与共同资源约束结合。

如果预算限制为0.80千元，B、C的总成本0.90超过预算，原最优方案不再可行。
加入约束 \(\sum_i c_i z_i\le0.80\)，枚举可得选择 B 单独加急的收益1.40，
高于 A、C 组合的1.30；B、D也为1.40，但多占一个名额且多耗0.20预算。
在收益并列时，可以将较少成本或较少干预规定为第二层选择标准。本例据此选择仅加急 B。

上述加法目标假定各订单收益可独立相加。如果加急 B 会挤占 C 的加工位置，
\(\delta_i\) 就可能依赖整个行动向量 \(z\)，问题必须改写为组合结果模型。
效果估计存在误差时，应围绕各订单的收益变化开展敏感性分析，判断哪些参数变化会改变最优组合，并报告相应的备选方案。

\section{人工智能辅助优化建模}
\subjectindex{qiujieqi}{求解器}
LLM 可辅助将自然语言约束转成候选模型、列出缺失参数和解释求解结果。可行性与最优性则应由求解器及独立检查负责。一个模型声称“所有任务均按时完成”，必须用实际开始、完成和交期逐项核验。

优化请求的执行流程包括：结构化输入模式验证，单位和范围检查，生成模型，求解，重算约束残差与目标，确认授权，再提交执行。求解器应报告可行性状态、时间限制、当前目标及有效界；达到时间限制时，应同时报告当前可行解和最优性间隙\footnote{最优性间隙衡量当前可行目标值与求解器有效界之间的差距。只有在界有效且间隙满足规定容差时，才能据此评价解距最优值有多远；相对间隙的分母约定应随报告说明。}。不可行时应返回相互冲突的约束与需要管理决策的放宽选项。

本地模型可通过预定义工具调用求解服务。服务限定输入参数与求解范围；涉及交期或质量约束的调整须由有权人员批准。第~\ref{ch:week10-process-mining}~章会把这些步骤放入可恢复的执行状态机。

\section{推理服务的部署与资源配置}
\subjectindex{bendiduanfuwu}{本地推理服务}
企业选用模型服务还面对一个资源配置问题。云端按调用量或 token 计价，本地
部署承担较大固定投入及持续运维；两者的数据流向、可用地区和并发能力也可能
不同。设每小时预计有 \(q\) 个同质请求，云端每请求全成本 \(c_c\)，本地服务
每小时固定成本 \(F\)、每请求可变成本 \(c_l\)。在服务质量、数据权限和
时延约束均满足的条件下，两项小时成本分别为
\[
  C_c(q)=qc_c,\qquad C_l(q)=F+qc_l.
\]
若 \(c_c>c_l\)，两个函数相交于 \(q^*=F/(c_c-c_l)\)。这个阈值只比较
所定义的费用，不说明本地方案在 \(q>q^*\) 时必然可用。还须检查本地容量
上限、到达时段的集中程度、维护停机和响应时间分布。

令本地系统有 \(k\) 个等效并行服务位置，每次调用平均占用 \(s\) 秒。若只把
\(\rho_{\max}\) 作为设计阶段的平均占用上限，则候选方案需满足
\[
  q\,s\le 3600k\rho_{\max}.
\]
该条件只约束平均占用，不能保证第95百分位完成时间满足要求；后者需要根据实际到达和
服务时间分布测量或仿真。权限与数据驻留要求是额外的可行性约束，不能通过
低价格抵消。

例9.8（模型服务的部署区间）。取教学假设 \(c_c=0.0029\) 美元/请求，
\(c_l=0.0002\) 美元/请求，\(F=3\) 美元/小时，\(s=6.186\) 秒，
\(k=4\)，\(\rho_{\max}=0.7\)。则费用交点约为1111请求/小时，
平均占用限制约为1629请求/小时。900请求/小时下，云端费用为2.61美元，
本地为3.18美元；1500请求/小时下，分别为4.35美元与3.30美元。
若超过本地设计容量，可增加实例、分流或改变模型，每种变化都会修改成本和
性能。此处的价格和速度均为教学设定，与任何厂商报价无关。附录
\ref{app:ai-infra-lab}~的程序逐项复算这些数字。

实际比较还需计入失败重试、人工检查、训练与版本管理、网络传输、数据存储及
服务中断的损失。多种请求共享设备时，前述同质请求假设也应改写，分别测量
短问答、长文档和多工具任务的服务时间。仅用一个平均值确定全部资源需求，
会掩盖日内峰值及长任务对队列的影响。

\section{数值求解与库存计算实例}
\subsection{求解方法}

小规模离散优化可以枚举所有可行方案。若每个产品只有几十个候选订货量，枚举法可逐项核对目标值与约束。大规模企业问题通常需要线性规划、混合整数规划、动态规划、启发式算法或商业求解器。

求解方法的选择取决于变量类型、约束结构、问题规模和计算期限。配套库存算例使用 Python 枚举，便于核对每个候选方案的目标值和资源消耗。

\subsection{库存配置计算}

本章样例数据位于：
\begin{center}
\path{data/sample/replenishment_planning.csv}
\end{center}
数据粒度为一行一个 SKU，包含单位采购成本、按前述 \(m\) 定义的单位销售收入、持有成本、缺货惩罚、当前库存、最大订货量、库容占用，以及低、中、高三种需求情景和概率。

配套代码位于：
\begin{center}
\path{src/bdm_decision/cases/week09_optimization.py}
\end{center}
核心优化工具位于：
\begin{center}
\path{src/bdm_decision/optimization/inventory_optimization.py}
\end{center}

在预算上限 2600、库容上限 95 的基准方案下，最优补货计划为：
\begin{table}[htbp]
\centering
\caption{库存枚举优化的基准方案；期望数量按需求情景加权}
\begin{tabular}{lrrrr}
\toprule
\rowcolor{BDMTableHead}
SKU & 订货量 & 期望销售 & 期望缺货 & 期望净价值\\
\midrule
P01 AirFilter & 18 & 28.00 & 2.00 & 1102.00\\
P02 SmartValve & 16 & 22.40 & 2.00 & 1113.60\\
P03 SensorKit & 10 & 27.40 & 6.70 & 877.50\\
P04 ControlBox & 8 & 12.25 & 4.00 & 821.50\\
\bottomrule
\end{tabular}
\end{table}
该方案使用预算 2596，使用库容 93，总期望净价值为 3914.60。

脚本还输出各产品前五个订货单位的边际价值：
\begin{table}[htbp]
\centering
\caption{各产品前五个新增订货单位的期望边际价值}
\begin{tabular}{lr}
\toprule
\rowcolor{BDMTableHead}
产品 & 前五个边际价值\\
\midrule
P01 AirFilter & 44.00, 44.00, 44.00, 44.00, 44.00\\
P02 SmartValve & 65.00, 65.00, 65.00, 65.00, 65.00\\
P03 SensorKit & 33.00, 33.00, 33.00, 33.00, 33.00\\
P04 ControlBox & 86.00, 86.00, 86.00, 86.00, 86.00\\
\bottomrule
\end{tabular}
\end{table}
P04 的早期边际价值最高，但它单位采购成本和库容占用也较高，因此基准方案不会仅按边际价值优先补充 P04。单品边际价值排序不能替代预算和库容约束下的组合比较。

\section{优化结果的实施与评价}

补货方案实施前，需要核对当前数据和执行条件，具体包括以下事项。
\begin{enumerate}
  \item 需求预测是否使用了最新促销、季节和客户订单信息；
  \item 单位成本、贡献毛利和缺货惩罚是否由业务部门确认；
  \item 供应商提前期是否允许按建议数量采购；
  \item 是否存在最小起订量、包装倍数或合同约束；
  \item SKU 之间是否有替代或互补关系；
  \item 方案是否对预算、库容和需求情景过度敏感；
  \item 执行后是否有反馈数据用于更新预测和优化参数。
\end{enumerate}

优化参数随执行反馈更新，形成以下运行过程：
\[
\text{预测}
\rightarrow
\text{目标与约束}
\rightarrow
\text{求解}
\rightarrow
\text{执行}
\rightarrow
\text{反馈更新}.
\]
环境变化后，输入参数需要根据反馈更新；未经核查的业务参数可能使求解结果偏离实际目标。

\subsection{模型运行与管理责任}

优化模型的治理规定目标、参数和约束的确认责任，以及例外处理和执行监测的职责。若建议减少某类产品库存，分析报告应说明需求情景、成本参数和库容限制各自的作用，并给出相关敏感性结果。

优化报告至少应包含：模型目标、主要约束、输入数据日期、关键参数来源、最优方案、资源使用、敏感性结果、模型未纳入的因素、执行建议和监控指标。对于高风险场景，还应保留人工复核机制和紧急回退方案。

\section*{本章小结}
\addcontentsline{toc}{section}{本章小结}

预测模型估计未来状态，因果推断评估行动效果，优化模型在约束下选择行动。一个规范的优化问题必须明确决策变量、目标函数、约束、可行域和输入假设。库存补货模型用离散需求情景计算期望净价值，并在预算和库容约束下选择订货量。方案实施前应进行敏感性分析和执行条件检查，实施后根据反馈更新参数。

线性规划对偶给出最优性证明和局部资源价格，报童模型把过量与短缺成本转化为订货分位数，CVaR 衡量尾部损失，整数规划处理不可分选择与活动顺序。决策后悔值衡量使用估计参数相对于完全信息最优方案的损失。模型服务的部署选择同样需要比较固定成本、单位费用和容量限制；费用交点仅确定经济候选范围。方案的实际执行及其经营结果可通过事件日志研究，见第10章。

\chapterreading{线性优化的理论与计算基础见 Dantzig 及 Bertsimas 与 Tsitsiklis
\cite{dantzig1963,bertsimas1997}，整数约束见 Nemhauser 与 Wolsey
\cite{nemhauser1988}，凸优化见 Boyd 与 Vandenberghe\cite{boyd2004}。
库存和随机规划可分别参考 Silver 等与 Shapiro 等\cite{silver2016,shapiro2021}。
从预测走向处方的研究见 Bertsimas 与 Kallus，决策导向学习见 Elmachtoub 与
Grigas\cite{bk2020,spo09}。学习和优化可以联合设计，其目标函数与现实损失之间的对应关系仍需独立论证。}

\section*{思考与习题}
\addcontentsline{toc}{section}{思考与习题}

\begin{exercise}[证书与成本]
验证线性规划例子的原、对偶解和目标值。将第一个资源上限从 4 调至 4.1，求新解并检查局部影子价格。说明为何整数变量时不能直接使用同一微分解释。
\end{exercise}
\begin{exercise}[排程与风险]
枚举两任务排程，重算加权延期。再构造两个期望损失相同而尾部损失不同的情景分布，比较 \(90\%\) CVaR。明确所有数字均为教学假设。
\end{exercise}

\begin{exercise}
围绕 \path{data/sample/replenishment_planning.csv} 完成下列任务：
\begin{enumerate}
  \item 写出本案例的决策变量、目标函数和两个主要约束；
  \item 对 P02 在订货量 16 时计算期望销售、期望缺货和期望净价值；
  \item 计算 P04 前五个订货单位的边际价值，并解释其管理含义；
  \item 将预算上限改为 2200，记录新最优方案；
  \item 将库容上限改为 75，记录新最优方案；
  \item 写一段不超过 400 字的管理建议，说明是否应直接执行基准方案，以及执行前还应验证哪些业务条件。
\end{enumerate}
\end{exercise}

\begin{exercise}[加急组合与边界条件]
对表~\ref{tab:expedite-choice}~列出全部至多选择两张订单的组合，分别求无预算限制、
预算0.80和预算0.90时的最优方案，保留并列解。再令 C 的效果从0.20降为未知数
\(d\)，求在无预算限制时 B、C 优于 A、B 的条件。讨论这一条件如何帮助决定是否
值得追加一轮效果试验。
\end{exercise}

\begin{exercise}[由运营模型反推产能]
从式~\eqref{eq:queue-capacity}~解释到达率、持有成本和能力成本如何改变最优能力
与利用率。若现实设备只能提供服务率9、10、12三个选项，应如何使用连续解帮助
比较而不把它直接当作可执行方案？若到达率本身不确定，说明还需要哪些情景和
风险判断。所有推断必须保留 M/M/1 的适用条件。
\end{exercise}
